% Options for packages loaded elsewhere
\PassOptionsToPackage{unicode}{hyperref}
\PassOptionsToPackage{hyphens}{url}
\documentclass[
]{article}
\usepackage{xcolor}
\usepackage[margin=1in]{geometry}
\usepackage{amsmath,amssymb}
\setcounter{secnumdepth}{-\maxdimen} % remove section numbering
\usepackage{iftex}
\ifPDFTeX
  \usepackage[T1]{fontenc}
  \usepackage[utf8]{inputenc}
  \usepackage{textcomp} % provide euro and other symbols
\else % if luatex or xetex
  \usepackage{unicode-math} % this also loads fontspec
  \defaultfontfeatures{Scale=MatchLowercase}
  \defaultfontfeatures[\rmfamily]{Ligatures=TeX,Scale=1}
\fi
\usepackage{lmodern}
\ifPDFTeX\else
  % xetex/luatex font selection
\fi
% Use upquote if available, for straight quotes in verbatim environments
\IfFileExists{upquote.sty}{\usepackage{upquote}}{}
\IfFileExists{microtype.sty}{% use microtype if available
  \usepackage[]{microtype}
  \UseMicrotypeSet[protrusion]{basicmath} % disable protrusion for tt fonts
}{}
\makeatletter
\@ifundefined{KOMAClassName}{% if non-KOMA class
  \IfFileExists{parskip.sty}{%
    \usepackage{parskip}
  }{% else
    \setlength{\parindent}{0pt}
    \setlength{\parskip}{6pt plus 2pt minus 1pt}}
}{% if KOMA class
  \KOMAoptions{parskip=half}}
\makeatother
\usepackage{color}
\usepackage{fancyvrb}
\newcommand{\VerbBar}{|}
\newcommand{\VERB}{\Verb[commandchars=\\\{\}]}
\DefineVerbatimEnvironment{Highlighting}{Verbatim}{commandchars=\\\{\}}
% Add ',fontsize=\small' for more characters per line
\usepackage{framed}
\definecolor{shadecolor}{RGB}{248,248,248}
\newenvironment{Shaded}{\begin{snugshade}}{\end{snugshade}}
\newcommand{\AlertTok}[1]{\textcolor[rgb]{0.94,0.16,0.16}{#1}}
\newcommand{\AnnotationTok}[1]{\textcolor[rgb]{0.56,0.35,0.01}{\textbf{\textit{#1}}}}
\newcommand{\AttributeTok}[1]{\textcolor[rgb]{0.13,0.29,0.53}{#1}}
\newcommand{\BaseNTok}[1]{\textcolor[rgb]{0.00,0.00,0.81}{#1}}
\newcommand{\BuiltInTok}[1]{#1}
\newcommand{\CharTok}[1]{\textcolor[rgb]{0.31,0.60,0.02}{#1}}
\newcommand{\CommentTok}[1]{\textcolor[rgb]{0.56,0.35,0.01}{\textit{#1}}}
\newcommand{\CommentVarTok}[1]{\textcolor[rgb]{0.56,0.35,0.01}{\textbf{\textit{#1}}}}
\newcommand{\ConstantTok}[1]{\textcolor[rgb]{0.56,0.35,0.01}{#1}}
\newcommand{\ControlFlowTok}[1]{\textcolor[rgb]{0.13,0.29,0.53}{\textbf{#1}}}
\newcommand{\DataTypeTok}[1]{\textcolor[rgb]{0.13,0.29,0.53}{#1}}
\newcommand{\DecValTok}[1]{\textcolor[rgb]{0.00,0.00,0.81}{#1}}
\newcommand{\DocumentationTok}[1]{\textcolor[rgb]{0.56,0.35,0.01}{\textbf{\textit{#1}}}}
\newcommand{\ErrorTok}[1]{\textcolor[rgb]{0.64,0.00,0.00}{\textbf{#1}}}
\newcommand{\ExtensionTok}[1]{#1}
\newcommand{\FloatTok}[1]{\textcolor[rgb]{0.00,0.00,0.81}{#1}}
\newcommand{\FunctionTok}[1]{\textcolor[rgb]{0.13,0.29,0.53}{\textbf{#1}}}
\newcommand{\ImportTok}[1]{#1}
\newcommand{\InformationTok}[1]{\textcolor[rgb]{0.56,0.35,0.01}{\textbf{\textit{#1}}}}
\newcommand{\KeywordTok}[1]{\textcolor[rgb]{0.13,0.29,0.53}{\textbf{#1}}}
\newcommand{\NormalTok}[1]{#1}
\newcommand{\OperatorTok}[1]{\textcolor[rgb]{0.81,0.36,0.00}{\textbf{#1}}}
\newcommand{\OtherTok}[1]{\textcolor[rgb]{0.56,0.35,0.01}{#1}}
\newcommand{\PreprocessorTok}[1]{\textcolor[rgb]{0.56,0.35,0.01}{\textit{#1}}}
\newcommand{\RegionMarkerTok}[1]{#1}
\newcommand{\SpecialCharTok}[1]{\textcolor[rgb]{0.81,0.36,0.00}{\textbf{#1}}}
\newcommand{\SpecialStringTok}[1]{\textcolor[rgb]{0.31,0.60,0.02}{#1}}
\newcommand{\StringTok}[1]{\textcolor[rgb]{0.31,0.60,0.02}{#1}}
\newcommand{\VariableTok}[1]{\textcolor[rgb]{0.00,0.00,0.00}{#1}}
\newcommand{\VerbatimStringTok}[1]{\textcolor[rgb]{0.31,0.60,0.02}{#1}}
\newcommand{\WarningTok}[1]{\textcolor[rgb]{0.56,0.35,0.01}{\textbf{\textit{#1}}}}
\usepackage{graphicx}
\makeatletter
\newsavebox\pandoc@box
\newcommand*\pandocbounded[1]{% scales image to fit in text height/width
  \sbox\pandoc@box{#1}%
  \Gscale@div\@tempa{\textheight}{\dimexpr\ht\pandoc@box+\dp\pandoc@box\relax}%
  \Gscale@div\@tempb{\linewidth}{\wd\pandoc@box}%
  \ifdim\@tempb\p@<\@tempa\p@\let\@tempa\@tempb\fi% select the smaller of both
  \ifdim\@tempa\p@<\p@\scalebox{\@tempa}{\usebox\pandoc@box}%
  \else\usebox{\pandoc@box}%
  \fi%
}
% Set default figure placement to htbp
\def\fps@figure{htbp}
\makeatother
\setlength{\emergencystretch}{3em} % prevent overfull lines
\providecommand{\tightlist}{%
  \setlength{\itemsep}{0pt}\setlength{\parskip}{0pt}}
\usepackage{bookmark}
\IfFileExists{xurl.sty}{\usepackage{xurl}}{} % add URL line breaks if available
\urlstyle{same}
\hypersetup{
  pdftitle={The Bellbeat Case Study: How can a wellness company play it smart?},
  pdfauthor={Yash Tiwari},
  hidelinks,
  pdfcreator={LaTeX via pandoc}}

\title{The Bellbeat Case Study: How can a wellness company play it
smart?}
\author{Yash Tiwari}
\date{2026-07-02}

\begin{document}
\maketitle

\subsection{Phase 1: Ask}\label{phase-1-ask}

\subsubsection{Business Task}\label{business-task}

To identify consumer trends in smart device usage from non-Bellabeat
users (Fitbit data) to help optimize and guide Bellabeat's marketing and
product development strategy.

\subsubsection{Key Stakeholders}\label{key-stakeholders}

\begin{itemize}
\tightlist
\item
  \textbf{Urška Sršen} -- Chief Creative Officer \& Co-founder
\item
  \textbf{Sando Mur} -- Mathematician \& Co-founder
\item
  \textbf{Bellabeat Marketing Analytics Team}
\end{itemize}

\begin{center}\rule{0.5\linewidth}{0.5pt}\end{center}

\subsection{Phase 2: Prepare}\label{phase-2-prepare}

\subsubsection{Data Source:}\label{data-source}

Data was gathered from the public dataset:
\href{https://www.kaggle.com/datasets/arashnic/fitbit}{Fitbit Tracker
Data via Kaggle}. This open data contains tracking logs from 30 eligible
Fitbit users who consented to submit personal tracker data, including
minute-level physical activity, heart rate, and sleep metrics.

\subsubsection{Critical Data
Limitations}\label{critical-data-limitations}

⚠️ \textbf{Important Analytical Disconnect:} The Fitbit dataset tracks
30 random users with \textbf{no provided gender demographics}. Because
Bellabeat specifically engineers health-focused smart products for
women, treating these calorie burns or step distances as a perfect 1:1
behavioral blueprint presents fundamental limitations. On average, men
exhibit higher basal metabolic rates (BMR) and longer average stride
lengths.

\emph{Conclusion:} While this public dataset serves as a useful proxy
for general human device interaction and user adherence, our strategic
findings must be validated against native, internal Bellabeat female
user metrics prior to executing high-cost marketing workflows.

\begin{center}\rule{0.5\linewidth}{0.5pt}\end{center}

\subsection{Phase 3: Process}\label{phase-3-process}

\subsubsection{Data Cleaning Log}\label{data-cleaning-log}

\begin{enumerate}
\def\labelenumi{\arabic{enumi}.}
\tightlist
\item
  \textbf{\texttt{dailyActivity\_merged} Verification:} Checked for
  structural anomalies. Zero duplicate rows found. Standardized
  character dates into true Date-formatted entities. Flagged tracking
  days where recorded logged minutes fell short of a full 1440-minute
  window. Investigated step count outliers.
\item
  \textbf{\texttt{minuteSleep\_merged} Sanitation:} Located and purged
  \textbf{525 duplicate records}. Isolated clean calendar dates from
  combined timestamp variables. Verified structural integrity and blank
  entries. Confirmed data capture over 23 unique user IDs.
\end{enumerate}

\subsubsection{Setting Up the
Environment}\label{setting-up-the-environment}

\begin{Shaded}
\begin{Highlighting}[]
\FunctionTok{library}\NormalTok{(tidyverse)}
\end{Highlighting}
\end{Shaded}

\begin{verbatim}
## Warning: package 'ggplot2' was built under R version 4.6.1
\end{verbatim}

\begin{verbatim}
## Warning: package 'readr' was built under R version 4.6.1
\end{verbatim}

\begin{verbatim}
## -- Attaching core tidyverse packages ------------------------ tidyverse 2.0.0 --
## v dplyr     1.2.1     v readr     2.2.0
## v forcats   1.0.1     v stringr   1.6.0
## v ggplot2   4.0.3     v tibble    3.3.1
## v lubridate 1.9.5     v tidyr     1.3.2
## v purrr     1.2.2     
## -- Conflicts ------------------------------------------ tidyverse_conflicts() --
## x dplyr::filter() masks stats::filter()
## x dplyr::lag()    masks stats::lag()
## i Use the conflicted package (<http://conflicted.r-lib.org/>) to force all conflicts to become errors
\end{verbatim}

\begin{Shaded}
\begin{Highlighting}[]
\FunctionTok{library}\NormalTok{(readr)}
\FunctionTok{library}\NormalTok{(janitor)}
\end{Highlighting}
\end{Shaded}

\begin{verbatim}
## Warning: package 'janitor' was built under R version 4.6.1
\end{verbatim}

\begin{verbatim}
## 
## Attaching package: 'janitor'
## 
## The following objects are masked from 'package:stats':
## 
##     chisq.test, fisher.test
\end{verbatim}

\begin{Shaded}
\begin{Highlighting}[]
\CommentTok{\# Load local cleaned datasets}
\NormalTok{daily\_activity }\OtherTok{\textless{}{-}} \FunctionTok{read.csv}\NormalTok{(}\StringTok{"02\_Cleaned\_data/dailyActivity\_merged\_clean.csv"}\NormalTok{)}
\NormalTok{daily\_sleep }\OtherTok{\textless{}{-}} \FunctionTok{read.csv}\NormalTok{(}\StringTok{"02\_Cleaned\_data/dailySleep\_Summary.csv"}\NormalTok{)}
\end{Highlighting}
\end{Shaded}

\subsubsection{Data Merging}\label{data-merging}

To evaluate dependencies between daytime habits and subsequent sleep
duration, datasets were combined via an inner join. This resulted in an
intersection of 199 completed observations containing concurrent daily
movement and sleep logs.

\begin{Shaded}
\begin{Highlighting}[]
\CommentTok{\# Column name standardization and merging}
\NormalTok{daily\_activity\_clean }\OtherTok{\textless{}{-}}\NormalTok{ daily\_activity }\SpecialCharTok{\%\textgreater{}\%} \FunctionTok{clean\_names}\NormalTok{()}
\NormalTok{daily\_sleep\_clean }\OtherTok{\textless{}{-}}\NormalTok{ daily\_sleep }\SpecialCharTok{\%\textgreater{}\%} \FunctionTok{clean\_names}\NormalTok{()}

\CommentTok{\# Merging \& Saving}

\NormalTok{merged\_data }\OtherTok{\textless{}{-}}\NormalTok{ daily\_activity\_clean }\SpecialCharTok{\%\textgreater{}\%} \FunctionTok{inner\_join}\NormalTok{(daily\_sleep\_clean, }\AttributeTok{by =} \FunctionTok{c}\NormalTok{(}\StringTok{"id"} \OtherTok{=} \StringTok{"id"}\NormalTok{, }\StringTok{"activity\_date"} \OtherTok{=} \StringTok{"sleep\_date"}\NormalTok{ ))}
\FunctionTok{write.csv}\NormalTok{(merged\_data, }\StringTok{"02\_Cleaned\_data/Daily\_activity\_sleeep\_merged.csv"}\NormalTok{)}
\end{Highlighting}
\end{Shaded}

\begin{center}\rule{0.5\linewidth}{0.5pt}\end{center}

\subsection{Phase 4: Analyze \&
Visualize}\label{phase-4-analyze-visualize}

\subsubsection{Daily Activity Profile
Breakdown}\label{daily-activity-profile-breakdown}

We analyzed consumer intensity distribution baselines across tracking
records to map active vs.~sedentary lifestyles:

\begin{Shaded}
\begin{Highlighting}[]
\NormalTok{activity\_summary\_table }\OtherTok{\textless{}{-}}\NormalTok{ daily\_activity }\SpecialCharTok{\%\textgreater{}\%} \FunctionTok{summarise}\NormalTok{(}\AttributeTok{average\_very\_active\_minutes =} \FunctionTok{mean}\NormalTok{(VeryActiveMinutes), }\AttributeTok{average\_fairly\_active\_minutes =} \FunctionTok{mean}\NormalTok{(FairlyActiveMinutes), }\AttributeTok{average\_lightly\_active\_minutes =} \FunctionTok{mean}\NormalTok{(LightlyActiveMinutes), }\AttributeTok{average\_sedantry\_minutes =} \FunctionTok{mean}\NormalTok{(SedentaryMinutes))}
\end{Highlighting}
\end{Shaded}


\end{document}
